\documentclass[aps,prl,twocolumn,superscriptaddress]{revtex4-2}
\usepackage{amsmath,amssymb,amsfonts}
\usepackage{graphicx}
\usepackage{hyperref}

\newcommand{\lean}[1]{\texttt{#1}}
\newcommand{\nativedecide}{\texttt{native\_decide}}

\begin{document}

\title{From Surgery to Invariants: The First Formalization of the\\
WRT TQFT Pipeline in Lean~4}

\author{J.~Roehm}
\affiliation{Independent Researcher}

\date{\today}

\begin{abstract}
We present the first formalization of the Witten-Reshetikhin-Turaev
3-manifold invariant pipeline in any proof assistant. Starting from
verified modular tensor category data (Ising and Fibonacci MTCs with
machine-checked R-matrices, F-symbols, pentagon, hexagon, and ribbon
conditions), we build a complete chain through the Temperley-Lieb
algebra, Jones-Wenzl idempotents, surgery presentations for
3-manifolds, to the WRT surgery formula connecting topological data
to algebraic invariants. Concrete computations verified by
\nativedecide{}: Ising $D^2 = 4$, Gauss sum $p_+ = 2\zeta_{16}$;
Fibonacci $D^2 = 2 + \varphi$ with $\varphi^2 = \varphi + 1$. The
Temperley-Lieb algebra (3~relation types), Jones-Wenzl recurrence,
surgery presentations ($S^3$, $S^2 \times S^1$, lens spaces), and the
WRT formula are all first formalizations in any proof assistant. Zero
sorry, zero axioms across 7~new Lean~4 modules.
\end{abstract}

\maketitle

%% ══════════════════════════════════════════════════════════════════
\section{Introduction}

The Witten-Reshetikhin-Turaev invariant~\cite{witten1989,rt1991}
assigns a complex number $Z(M)$ to each closed oriented 3-manifold $M$,
computed from a modular tensor category (MTC). Despite 35~years of
development and central importance in quantum topology, no component
of the WRT construction has previously been formalized in any proof
assistant --- not the surgery formula, not the Temperley-Lieb algebra,
not the Kirby calculus, and not any quantum 3-manifold invariant value.

We close this gap with a complete formalization pipeline in Lean~4
with Mathlib, building on our prior work establishing the first
verified MTC instances (Ising and Fibonacci) with machine-checked
pentagon, hexagon, and ribbon coherence~\cite{paper14}. The pipeline
comprises seven new modules totaling ${\sim}60$~theorems with zero
\texttt{sorry} and zero axioms.

%% ══════════════════════════════════════════════════════════════════
\section{The Temperley-Lieb Algebra}

The Temperley-Lieb algebra $\mathrm{TL}_n(\delta)$ is the algebraic
foundation of the BHMV construction~\cite{bhmv1995} of the WRT TQFT
functor. We define it in \lean{TemperleyLieb.lean} as a quotient of
the free algebra on $n$ generators $\{e_0, \ldots, e_{n-1}\}$ by
three relation families:
\begin{align}
  \text{(TL1)} && e_i^2 &= \delta \cdot e_i \label{eq:tl1} \\
  \text{(TL2)} && e_i e_j e_i &= e_i \quad (|i-j|=1) \label{eq:tl2} \\
  \text{(TL3)} && e_i e_j &= e_j e_i \quad (|i-j| \ge 2) \label{eq:tl3}
\end{align}
All three relations are proved via the standard
\lean{RingQuot.mkAlgHom\_rel} descent from the free algebra. The
parameter $\delta$ is the loop value; for SU(2)$_k$ quantum groups,
$\delta = [2]_q = q + q^{-1}$.

%% ══════════════════════════════════════════════════════════════════
\section{Jones-Wenzl Idempotents}

The Jones-Wenzl idempotent $f^{(m)} \in \mathrm{TL}_n$ is defined by
the Wenzl recurrence~\cite{wenzl1988}:
\begin{equation}
  f^{(0)} = 1, \quad
  f^{(m+1)} = f^{(m)} - c(m) \cdot f^{(m)} \cdot e_m \cdot f^{(m)}
\end{equation}
where $c(m) = [m]_q / [m+1]_q$ for the quantum-group specialization.
We parametrize by an abstract coefficient function $c : \mathbb{N} \to k$,
keeping the definition valid over any commutative ring. The recurrence
step, base case $f^{(0)} = 1$, and explicit $f^{(1)} = 1 - c(0) \cdot e_0$
are all definitional equalities (proved by \lean{rfl}).

%% ══════════════════════════════════════════════════════════════════
\section{Surgery Presentations}

Every closed oriented 3-manifold is obtained by Dehn surgery on a
framed link in $S^3$ (Lickorish-Wallace theorem). We encode surgery
presentations as symmetric integer linking matrices
(\lean{SurgeryPresentation.lean}), with standard examples:
\begin{itemize}
  \item $S^3$: empty surgery (0 components)
  \item $S^2 \times S^1$: 0-framed unknot (1 component, framing $= 0$)
  \item $L(p,1)$: $p$-surgery on the unknot
  \item Hopf link: 2 components with linking number 1
\end{itemize}
Framing and linking properties verified for all examples. The key
insight from~\cite{turaev2010}: working with linking matrices avoids
smooth 3-manifold topology entirely --- the construction is purely
combinatorial.

%% ══════════════════════════════════════════════════════════════════
\section{The WRT Invariant Formula}

The WRT invariant connects MTC data to 3-manifold topology via the
surgery formula. We define \lean{MTCWithTwist} extending
\lean{PreModularData} with topological twist eigenvalues $\theta_i$,
global dimension $D^2 = \sum_i d_i^2$, and Gauss sum
$p_+ = \sum_i d_i^2 \theta_i$. The invariant for the empty surgery is:
\begin{equation}
  Z(S^3) = p_+ \cdot D^{-2}
\end{equation}
and $Z(S^2 \times S^1)$ equals the rank (number of simple objects),
proved as a definitional equality.

%% ══════════════════════════════════════════════════════════════════
\section{Verified Computations}

Using verified MTC data from companion work~\cite{paper14}:

\textbf{Ising MTC} (SU(2)$_2$, 3 simple objects $\{1, \sigma, \psi\}$):
\begin{itemize}
  \item $D^2 = 1 + (\sqrt{2})^2 + 1 = 4$ \quad (\nativedecide{} over $\mathbb{Q}(\zeta_{16})$)
  \item $p_+ = 1 + 2\theta_\sigma + \theta_\psi = 2\zeta_{16}$ \quad (proved)
  \item $Z(S^2 \times S^1) = 3$ \quad (rank)
  \item Braiding gates: $\sigma_1 = \mathrm{diag}(\zeta^{-1}, \zeta^3)$ (Clifford)
\end{itemize}

\textbf{Fibonacci MTC} (2 simple objects $\{1, \tau\}$, $\tau \otimes \tau = 1 \oplus \tau$):
\begin{itemize}
  \item $\varphi^2 = \varphi + 1$ \quad (\nativedecide{} over $\mathbb{Q}(\zeta_5)$)
  \item $D^2 = 1 + \varphi^2 = 2 + \varphi$ \quad (proved)
  \item $\theta_\tau = \zeta_5^2 = e^{4\pi i/5}$ \quad (proved)
  \item $Z(S^2 \times S^1) = 2$ \quad (rank)
\end{itemize}

%% ══════════════════════════════════════════════════════════════════
\section{q-Binomial Coefficients}

The q-binomial coefficient $\binom{n}{j}_q$ is defined via the
Laurent Pascal recurrence without division (\lean{QNumber.lean}):
\begin{equation}
  \binom{n{+}1}{j{+}1}_q = q^{j+1} \binom{n}{j{+}1}_q + q^{j-n} \binom{n}{j}_q
\end{equation}
with classical limit $\binom{n}{j}_1 = C(n,j)$ proved by induction.
Key values: $\binom{2}{1}_q = [2]_q$ (simply-laced Serre coefficient),
$\binom{3}{1}_q = [3]_q$ (for $B_n/C_n$ types).

%% ══════════════════════════════════════════════════════════════════
\section{Conclusion}

We have presented the first formalization of the WRT 3-manifold
invariant pipeline: from algebraic MTC data through the Temperley-Lieb
algebra, Jones-Wenzl idempotents, and surgery presentations to
concrete verified invariant values. Seven new Lean~4 modules,
${\sim}60$~theorems, zero sorry, zero axioms. The Ising and Fibonacci
computations demonstrate that verified MTC data can produce verified
topological invariants --- the first such connection in any proof
assistant.

The infrastructure opens direct paths to topological quantum
computation (braiding gates from the $B = F^{-1}RF$ formula),
verified quantum error correction (Levin-Wen codes from fusion data),
and the full WRT TQFT functor (Kirby invariance proofs).

\begin{acknowledgments}
Lean~4 proofs verified by \texttt{lake build} (v4.28.0,
Mathlib \texttt{8f9d9cff}). Automated proofs by the Aristotle
theorem prover (Harmonic).
\end{acknowledgments}

\begin{thebibliography}{99}
\bibitem{witten1989}
E.~Witten, Comm.\ Math.\ Phys.\ \textbf{121}, 351 (1989).

\bibitem{rt1991}
N.~Reshetikhin, V.~Turaev, Invent.\ Math.\ \textbf{103}, 547 (1991).

\bibitem{bhmv1995}
C.~Blanchet, N.~Habegger, G.~Masbaum, P.~Vogel, Topology \textbf{34}, 883 (1995).

\bibitem{wenzl1988}
H.~Wenzl, Invent.\ Math.\ \textbf{92}, 349 (1988).

\bibitem{turaev2010}
V.~G.~Turaev, \emph{Quantum Invariants of Knots and 3-Manifolds} (de Gruyter, 2010).

\bibitem{paper14}
J.~Roehm, companion paper: Braided modular tensor categories and knot invariants in Lean~4.

\bibitem{paper11}
J.~Roehm, companion paper: $U_q(\mathfrak{sl}_2)$ in Lean~4: the first quantum group in a proof assistant.
\end{thebibliography}

\end{document}
